% ---------------------------------------------------------
% Definition --- General Solution
% ---------------------------------------------------------

\begin{definitionbox}{Definition}
\begin{definition}[General Solution]
\label{def:ordinary-differential-equations-solutions-general-solution}
A \textbf{general solution} of a differential equation is a family of solutions containing one or more arbitrary constants.

For an $n$th-order ordinary differential equation, the general solution typically contains $n$ arbitrary constants:
\[
y(t)=\Phi(t,C_1,C_2,\dots,C_n).
\]
Each choice of the constants $C_1,C_2,\dots,C_n$ gives a particular solution of the differential equation.
\end{definition}
\end{definitionbox}

\begin{remark*}[Standard quantified statement]
This definition fixes the named kind of differential-equation object by the
conditions stated in the definition.
\end{remark*}

\NoLocalDependencies

\begin{remark*}[Interpretation]
The general solution represents the full family of solutions produced by solving the differential equation. Initial conditions are then used to choose specific values of the arbitrary constants, producing a particular solution.
\end{remark*}

\begin{example*}[General Solution]
The differential equation
\[
y'=3y
\]
has general solution
\[
y(t)=Ce^{3t},
\]
where $C$ is an arbitrary constant. If an initial condition such as
\[
y(0)=5
\]
is imposed, then $C=5$, and the particular solution is
\[
y(t)=5e^{3t}.
\]
\end{example*}



% ---------------------------------------------------------
% Definition --- Family of Solutions
% ---------------------------------------------------------

\begin{definitionbox}{Definition}
\begin{definition}[Family of Solutions]
\label{def:ordinary-differential-equations-solutions-family-of-solutions}
A \textbf{family of solutions} of a differential equation is a collection of solution functions described by one or more parameters.

For example, a family of solutions may be written in the form
\[
y(t)=\Phi(t,C),
\]
or more generally,
\[
y(t)=\Phi(t,C_1,C_2,\dots,C_n),
\]
where $C,C_1,\dots,C_n$ are parameters. Each choice of the parameters gives one particular solution.
\end{definition}
\end{definitionbox}

\begin{remark*}[Standard quantified statement]
This definition fixes the named kind of differential-equation object by the
conditions stated in the definition.
\end{remark*}

\NoLocalDependencies

\begin{remark*}[Interpretation]
A family of solutions records many solutions at once. The parameters act as labels that distinguish one solution curve from another. When initial conditions are supplied, they usually determine the parameter values and select a single member of the family.
\end{remark*}

\begin{example*}[A Family of Solutions]
The differential equation
\[
y'=3y
\]
has the family of solutions
\[
y(t)=Ce^{3t},
\]
where $C\in\mathbb{R}$. Each value of $C$ gives a different solution.
\end{example*}


% ---------------------------------------------------------
% Definition --- Initial Condition
% ---------------------------------------------------------

\begin{definitionbox}{Definition}
\begin{definition}[Initial Condition]
\label{def:ordinary-differential-equations-solutions-initial-condition}
An \textbf{initial condition} specifies the value of the unknown function at a particular value of the independent variable.

For a first-order differential equation involving an unknown function $y=y(t)$, an initial condition has the form
\[
y(t_0)=y_0,
\]
where $t_0$ is the initial time and $y_0$ is the initial value.
\end{definition}
\end{definitionbox}

\begin{remark*}[Standard quantified statement]
This definition fixes the named kind of differential-equation object by the
conditions stated in the definition.
\end{remark*}

\NoLocalDependencies

\begin{remark*}[Interpretation]
An initial condition selects one solution from a family of solutions. If a differential equation has general solution
\[
y(t)=\Phi(t,C),
\]
then the condition
\[
y(t_0)=y_0
\]
is used to determine the constant $C$.
\end{remark*}

\begin{example*}[Initial Condition]
The differential equation
\[
y'=3y
\]
has general solution
\[
y(t)=Ce^{3t}.
\]
If the initial condition is
\[
y(0)=5,
\]
then
\[
5=Ce^{0}=C.
\]
Thus the particular solution satisfying the initial condition is
\[
y(t)=5e^{3t}.
\]
\end{example*}


% ---------------------------------------------------------
% Definition --- Initial Value Problem
% ---------------------------------------------------------

\begin{definitionbox}{Definition}
\begin{definition}[Initial Value Problem]
\label{def:ordinary-differential-equations-solutions-initial-value-problem}
An \textbf{initial value problem}, or \textbf{IVP}, is a differential equation together with one or more initial conditions.

For a first-order differential equation, an initial value problem has the form
\[
y'=f(t,y),
\qquad
y(t_0)=y_0.
\]
The differential equation describes how the function changes, while the initial condition specifies the value of the solution at the initial time $t_0$.
\end{definition}
\end{definitionbox}

\begin{remark*}[Standard quantified statement]
This definition fixes the named kind of differential-equation object by the
conditions stated in the definition.
\end{remark*}

\NoLocalDependencies

\begin{remark*}[Interpretation]
An initial value problem asks for the particular solution of a differential equation that passes through a specified point. Geometrically, the condition
\[
y(t_0)=y_0
\]
means that the solution curve must pass through the point
\[
(t_0,y_0).
\]
\end{remark*}

\begin{example*}[Initial Value Problem]
Consider the initial value problem
\[
y'=3y,
\qquad
y(0)=5.
\]
The differential equation has general solution
\[
y(t)=Ce^{3t}.
\]
Using the initial condition,
\[
5=y(0)=Ce^0=C.
\]
Therefore the solution of the initial value problem is
\[
y(t)=5e^{3t}.
\]
\end{example*}


% ---------------------------------------------------------
% Definitions --- Existence, Uniqueness, and Interval of Existence
% ---------------------------------------------------------

\begin{definitionbox}{Definition}
\begin{definition}[Existence]
\label{def:ordinary-differential-equations-solutions-existence}
An initial value problem is said to have \textbf{existence} if there is at least one function that satisfies both the differential equation and the initial condition.

For a first-order initial value problem
\[
y'=f(t,y),
\qquad
y(t_0)=y_0,
\]
existence means that there is a function $y(t)$, defined on some interval containing $t_0$, such that
\[
y'(t)=f(t,y(t))
\]
and
\[
y(t_0)=y_0.
\]
\end{definition}
\end{definitionbox}

\begin{remark*}[Standard quantified statement]
This definition fixes the named kind of differential-equation object by the
conditions stated in the definition.
\end{remark*}

\NoLocalDependencies

\begin{definitionbox}{Definition}
\begin{definition}[Uniqueness]
\label{def:ordinary-differential-equations-solutions-uniqueness}
An initial value problem is said to have \textbf{uniqueness} if it has at most one solution on the interval under consideration.

For a first-order initial value problem
\[
y'=f(t,y),
\qquad
y(t_0)=y_0,
\]
uniqueness means that if $y_1(t)$ and $y_2(t)$ are both solutions on the same interval $I$ containing $t_0$, then
\[
y_1(t)=y_2(t)
\]
for every $t\in I$.
\end{definition}
\end{definitionbox}

\begin{remark*}[Standard quantified statement]
This definition fixes the named kind of differential-equation object by the
conditions stated in the definition.
\end{remark*}

\NoLocalDependencies

\begin{definitionbox}{Definition}
\begin{definition}[Interval of Existence]
\label{def:ordinary-differential-equations-solutions-interval-of-existence}
An \textbf{interval of existence} for a solution of an initial value problem is an interval on which the solution is defined and satisfies the differential equation.

For an initial value problem
\[
y'=f(t,y),
\qquad
y(t_0)=y_0,
\]
an interval $I$ is an interval of existence if $t_0\in I$ and there is a solution $y:I\to\mathbb{R}$ such that
\[
y'(t)=f(t,y(t))
\]
for every $t\in I$, and
\[
y(t_0)=y_0.
\]
\end{definition}
\end{definitionbox}

\begin{remark*}[Standard quantified statement]
This definition fixes the named kind of differential-equation object by the
conditions stated in the definition.
\end{remark*}

\NoLocalDependencies

\begin{remark*}[Interpretation]
Existence asks whether at least one solution passes through the initial point. Uniqueness asks whether only one solution passes through that point. The interval of existence records how long the solution is valid.
\end{remark*}

\begin{example*}[Existence and Uniqueness]
The initial value problem
\[
y'=3y,
\qquad
y(0)=5
\]
has the solution
\[
y(t)=5e^{3t}.
\]
This solution exists and is unique on $\mathbb{R}$.

The interval of existence is
\[
(-\infty,\infty).
\]
\end{example*}
